% =============================================================================
%  RL-H4-01 | Draft pendahuluan, related work, dan metode (TERISI PENUH)
% -----------------------------------------------------------------------------
%  Role    : RL — Research Lead (@alim)
%  Sprint  : H4 | PARALLEL (sisa waktu di luar RL-H4-02)
%  Ticket  : RL-H4-01                          Status: DRAFT COMPLETED (H4)
%  Target  : reports/sprint1/01_pendahuluan.md, 02_related_work.md, 03_metode.md
%            & reports/sprint1/RL-H4-01__report_sections.tex
%  Input   : brief v2 §1–§6, docs/positioning.md, docs/eval_protocol_v1.md, docs/lit/
%  Output  : draft lengkap 3 bagian (Pendahuluan, Related Work, Metode)
%  Handoff : RL-H7-01
%  DoD     : [x] semua klaim novelty memakai "sepengetahuan kami"
%            [x] tidak ada klaim kategori JANGAN DIKLAIM
%            [x] persamaan Modul 1 sama dengan brief §5
% =============================================================================
\documentclass[11pt,a4paper]{article}
\input{00_preamble}

\title{\jobid{RL-H4-01} | Draft Laporan Sprint 1\\
\large Pendahuluan, Related Work, dan Metode}
\author{RL (\texttt{@alim}) \\ \small RENARO Lab}
\date{Kamis, 26 November 2026 (H4) \quad | \quad \textcolor{blue!70!black}{DRAFT LENGKAP H4}}

\begin{document}
\maketitle

\section{Pendahuluan}
Harga komoditas hortikultura strategis di Indonesia, khususnya cabai merah keriting,
saling terhubung antarprovinsi melalui jaringan distribusi perdagangan antardaerah,
tetapi kekuatan dan jeda transmisinya tidak bersifat statis maupun serentak.
Sentra produksi cabai terkonsentrasi di sejumlah provinsi utama (seperti Jawa Timur,
Jawa Tengah, Jawa Barat, Sumatera Utara, dan Sulawesi Selatan), sedangkan konsumsi
tersebar di seluruh 34 provinsi. Perubahan pasokan di sentra pemasok memerlukan waktu
transmisi lintas hari hingga beberapa minggu sebelum tercermin pada harga pasar
tradisional di provinsi tujuan.

Ketika wilayah produksi asal mengalami kekeringan berkepanjangan (\textit{dry spell}),
hujan ekstrem 5-harian, atau pergeseran rezim makro ENSO (El Ni\~no dan La Ni\~na),
jadwal panen, kualitas pasokan, dan pilihan jalur distribusi ikut bergeser sehingga
bobot keterkaitan dan jeda transmisi harga berubah. Model peramalan konvensional yang
mengasumsikan struktur graf dan jeda tetap mengalami lonjakan galat justru pada rezim
iklim ekstrem, padahal periode tersebut paling krusial untuk peringatan dini harga pangan
pada horizon 5--20 hari kerja ke depan.

Di sisi literatur, model graf dinamis harga pangan mutakhir seperti RCSGNet \cite{rcsgnet2026}
dan MarketDNA \cite{marketdna2026} menyimpulkan rezim pasar secara laten dari histori harga
semata (\textit{price-only}) tanpa variabel iklim eksogen dan tanpa distribusi jeda per edge.
Sebaliknya, model graf berjeda seperti MillGNN \cite{millgnn2025} dan SAMSGL \cite{samsgl2023}
tidak mengondisikan jeda pada iklim teramati di wilayah asal dan tidak diuji pada pergeseran
rezim iklim OOD. Sementara itu, studi yang menyertakan cuaca pada GNN harga pertanian
(CGNN \cite{cgnn2023}; Min et al., 2025) memperlakukan cuaca sebatas fitur node lokal, dan
studi jeda cuaca wilayah sumber pada teh Sri Lanka \cite{tea2026} bekerja tanpa struktur graf.

Studi ini menguji empat pertanyaan penelitian melalui tiga \textit{gate} keputusan:
\textbf{RQ0} (kelayakan sinyal: apakah iklim wilayah asal berjeda meningkatkan akurasi
LightGBM dibanding harga saja dan harga + iklim lokal; telah lulus Gate 1 pada H3 dengan
signifikansi DM pada \texttt{el\_nino} $p=0{,}000117$ dan \texttt{O1 dry} $p=0{,}0104$);
\textbf{RQ1} (metode inti: apakah graf berjeda terkondisi iklim asal pada Modul 1 mengalahkan
baseline temporal, ST-GNN generik, lag-graph ala MillGNN, dan rezim laten ala RCSGNet pada
split/seed/backbone identik, serta memulihkan jeda pada data sintetis);
\textbf{RQ2} (apakah \textit{reliability gate} pada Modul 2--3 menekan \textit{worst-regime error}
tanpa merusak periode normal); dan
\textbf{RQ3} (apakah model tahan terhadap kesalahan spesifikasi prior graf BPS 10--50\%
yang terjadi bersamaan dengan pergeseran rezim iklim).

Kontribusi penelitian diurutkan dari \textbf{benchmark $\to$ bukti empiris $\to$ metode}
dengan kualifikasi eksplisit pada setiap klaim kebaruan:
\begin{itemize}
  \item \textbf{Benchmark (Kontribusi Utama):} ``Sepengetahuan kami, studi ini menyediakan
        benchmark evaluasi pertama untuk harga cabai merah keriting antarprovinsi Indonesia
        (34 provinsi, 2018--2026) yang terstratifikasi lima rezim iklim fisik dengan fitur
        \textit{availability-safe} (\texttt{available\_at} $\le t$), holdout tersegel,
        \textit{worst-regime error}, dan uji misspesifikasi graph prior.''
  \item \textbf{Bukti Empiris:} ``Sepengetahuan kami, kami menyediakan perbandingan
        \textit{head-to-head} terkontrol antara pengondisian iklim fisik teramati di wilayah
        asal melawan rezim pasar laten berbasis harga (ala RCSGNet) dan graf berjeda tanpa
        iklim (ala MillGNN), disertai pelaporan transparan atas kondisi di mana sinyal iklim
        tidak membantu.''
  \item \textbf{Metode:} ``Sepengetahuan kami, belum ada ST-GNN antarprovinsi untuk harga
        cabai Indonesia yang secara eksplisit mengondisikan distribusi probabilitas jeda
        per edge berarah ($\ell=1,\dots,8$ pentad + slot \texttt{null}) pada iklim fisik
        teramati di provinsi asal.''
\end{itemize}

Ruang lingkup Sprint 1 dikunci pada \textbf{Scope A}: satu komoditas utama (cabai merah
keriting pasar tradisional PIHPS, SP2KP sebagai validasi silang), 34 provinsi struktur lama,
horizon $h \in \{5, 10, 20\}$ hari kerja (utama $h=10$), 3 \textit{rolling-origin} + 1 holdout
tersegel ($\ge$ 2026-06-01), dan 5 seed (\texttt{0..4}). \textbf{Batas studi:} seluruh bobot
edge $A^{(\ell)}_{ij,t}$ diinterpretasikan secara ketat sebagai \textit{dependensi prediktif
berarah yang terkondisi iklim}, bukan estimasi arus pasokan fisik harian dan bukan identifikasi
kausal struktural. Klaim jeda dibatasi pada resolusi pentad karena data harian CHIRPS v3
adalah disagregasi pentad, dan data mentah PIHPS tidak diredistribusi.

\section{Related work}
Literatur terkait dikelompokkan ke dalam empat kategori sesuai Brief v2 \S2 dan
\fpath{docs/positioning.md}:

\paragraph{1. Pembelajaran graf spatio-temporal dan harga pertanian.}
Graph WaveNet dan MTGNN memelopori pembelajaran graf adaptif, tetapi matriks ketetanggaannya
bersifat statis sepanjang periode uji. Pada harga pertanian, RCSGNet \cite{rcsgnet2026}
dan MarketDNA \cite{marketdna2026} memodelkan graf dinamis berbasis state pasar laten dari
harga semata (\textit{price-only}) tanpa iklim eksogen dan tanpa dimensi jeda per edge.
Sementara itu, CGNN \cite{cgnn2023} dan Min et al.\ (2025) menyertakan variabel cuaca hanya
sebagai fitur node lokal di atas graf statis atau \textit{self-attention} standar.

\paragraph{2. Graf berarah yang peka terhadap jeda waktu (\textit{lag/delay-aware}).}
MillGNN \cite{millgnn2025} mempelajari relasi \textit{lead-lag} multi-skala dengan bias
peluruhan $-\lambda\ell$, dan SAMSGL \cite{samsgl2023} menyelaraskan deret berjeda via FFT;
keduanya menurunkan bobot jeda dari jendela sinyal input tanpa pengondisian iklim wilayah asal.
CRAFormer \cite{craformer2025} memodelkan DAG berjeda dengan curah hujan lokal, tetapi pada
domain geoteknik longsor satu situs dan bukan jaringan pasar pangan multiregional.

\paragraph{3. Ketahanan OOD, rezim, dan fusi graf statis-dinamis.}
DRAN \cite{dran2025} dan SDGF \cite{sdgf2025} menggunakan gerbang atensi untuk menggabungkan
relasi statis dan korelasi dinamis, sedangkan STEVE \cite{ji2025steve} dan DynaSTar
\cite{hao2026dynastar} memodelkan \textit{unseen confounder} laten pada OOD trafik perkotaan.
Karena mekanisme graf adaptif, \textit{gating} statis-dinamis, dan regularisasi invarian telah
mapan pada kelompok ini, studi kami tidak mengklaim mekanisme tersebut sebagai kebaruan umum.

\paragraph{4. Harga pangan, jeda cuaca wilayah sumber, dan guncangan ENSO.}
Studi harga teh Sri Lanka \cite{tea2026} menunjukkan bahwa jeda cuaca 1--3 minggu di wilayah
perkebunan sumber meningkatkan akurasi prediksi harga, tetapi dalam kerangka regresi deret
waktu tanpa struktur graf spasial antarwilayah.

\paragraph{Tiga kalimat posisi resmi.}
(1) \textit{Sepengetahuan kami}, belum ada ST-GNN antarprovinsi untuk harga pangan strategis
Indonesia yang mengondisikan distribusi jeda transmisi per edge berarah ($\ell=1,\dots,8$
pentad + slot \texttt{null}) pada indikator iklim fisik teramati di provinsi asal ($c_{j,\le t}$).
(2) \textit{Sepengetahuan kami}, studi ini menyajikan benchmark harga cabai antarprovinsi
pertama yang menggabungkan fitur \textit{availability-safe} (\texttt{available\_at} $\le t$),
stratifikasi lima rezim iklim fisik dengan \textit{Worst-Regime MAE}, perbandingan
\textit{head-to-head} melawan rezim laten berbasis harga, dan uji kesalahan prior graf 10--50\%.
(3) Kami tidak mengklaim graf adaptif, \textit{gating} statis-dinamis, regularisasi invarian,
atau penyertaan cuaca pada GNN pertanian sebagai kebaruan umum; bobot graf dimaknai murni
sebagai dependensi prediktif berarah, bukan arus pasokan fisik maupun kausalitas struktural.

\section{Metode}
\subsection{Representasi data dan aturan ketersediaan}
Node merepresentasikan $N=34$ provinsi struktur lama Indonesia; jendela input $W=60$ hari
kerja memuat harga historis \textit{as-of}, kalender SKB, curah hujan CHIRPS v3 terbobot
produksi/luas (wilayah sendiri dan asal), serta indeks RONI vintage. Mengacu pada
\fpath{docs/data_contract.md} dan \fpath{reports/qa/availability.csv}, setiap keluarga fitur
membawa stempel \texttt{available\_at}: harga menggunakan \texttt{next\_workday} (lag 1 hari
kerja), CHIRPS menggunakan Opsi B utama (\texttt{pentad\_end} + 2 hari; Opsi A sensitivitas
5 hari), RONI menggunakan \texttt{vintage\_or\_conservative\_3m}, bobot produksi menggunakan
\texttt{year\_plus\_1}, dan kalender bersifat \texttt{known\_in\_advance}. Hanya observasi
dengan \texttt{available\_at} $\le t$ yang dipakai pada \textit{decision date} $t$, diuji
otomatis oleh \fpath{tests/test_no_leakage.py}. Label rezim (\fpath{regime_labels.csv})
bersifat \textit{eval-only} dan tidak pernah masuk fitur maupun gerbang model.

\subsection{Modul 1: climate-conditioned lag graph}
Untuk tujuan $i$, asal $j$, dan bin jeda pentad $\ell \in \{1,\dots,L\}$ ($L=8$):
\begin{equation}
A^{(\ell)}_{ij,t}
= \operatorname{softmax}_{j,\ell}
\left[
\phi_\theta\left(h_{i,\le t},\, h_{j,\le t},\, c_{j,\le t},\, s_{ij},\, \ell\right)
\right],
\label{eq:module1}
\end{equation}
dengan $h_{i,\le t}$ state encoder temporal bersama (GRU 1 lapis, hidden $d_h=64$),
$e_{j,t}=\mathrm{MLP}_{\text{clim}}(c_{j,\le t})$ encoder iklim asal atas 8 pentad terakhir
yang tersedia pada $t$ (\texttt{rain\_pentad\_anom}, \texttt{dry\_spell\_days},
\texttt{extreme\_5d}, \texttt{rain\_30d\_anom}, \texttt{roni}), $s_{ij}$ fitur edge statis
(log jarak, bendera biner \texttt{bps\_mask}, pangsa BPS), dan $\mathrm{emb}(\ell)$
\textit{embedding} bin jeda. Representasi berjeda $z_j^{(\ell)}$ adalah rata-rata state
temporal $z_j(\tau)$ pada jendela pentad ke-$\ell$ ke belakang ($\tau \in [t-5\ell, t-5\ell+4]$).
Kandidat $j \in \mathcal{C}(i)$ dibatasi \texttt{candidate\_mask} (varian \texttt{bps}:
mask BPS $\cup$ $k$-tetangga terdekat; varian \texttt{none}: semua node lain) disertai satu
opsi slot \texttt{null} bervektor nol agar provinsi tujuan dapat mengabaikan tetangga saat
sinyal bising. Agregasi pesan dan head prediksi multi-horizon dirumuskan sebagai:
\begin{equation}
m_{i,t} = \sum_{j \in \mathcal{C}(i)} \sum_{\ell=1}^{L} A^{(\ell)}_{ij,t}\, W_v\, z_j^{(\ell)},
\qquad
\hat{y}_{i,t+h} = \mathrm{MLP}_{\text{head}}\!\left([h_{i,\le t},\, m_{i,t},\, \mathrm{cal}_{t+h}]\right),
\end{equation}
dilatih dengan loss MAE pada \texttt{y\_mask} dan regularisasi entropi perhatian ringan
$\lambda_H H(A)$. Untuk $N=34$ node dan $L=8$ bin pentad, ruang skor pasangan-jeda maksimal
hanya $34 \times 34 \times 8 = 9.248$ entri per sampel (memori GPU puncak $< 3\text{ GB}$
pada batch 32).

\subsection{Modul 2--3: stable/adaptive dan reliability gate}
Untuk menekan varians pada periode tenang sekaligus menjaga kepekaan pada rezim ekstrem:
\begin{equation}
z_{t,i} = g_{t,i} \odot H^{\mathrm{adapt}}_{t,i} + (1 - g_{t,i})\odot H^{\mathrm{stable}}_{t,i},
\qquad
g_{t,i} = \sigma\left(\mathrm{MLP}_g([\text{shock}_{t,i},\text{rel}_{t,i},\text{shift}_{t,i}])\right),
\end{equation}
dengan $H^{\mathrm{adapt}}_{t,i} = m_{i,t}$ (keluaran Modul 1) dan $H^{\mathrm{stable}}_{t,i}$
agregasi graf statis terpelajari berregularisasi $L_2$ kuat (atau prior BPS ternormalisasi).
Seluruh input gate tersedia pada $t$: $\text{shock}_{t,i}$ (agregat \texttt{extreme\_5d} dan
\texttt{dry\_spell\_days} asal serta $|\text{RONI}_t|$), $\text{rel}_{t,i}$ (entropi $A_t[i]$
dan jaraknya dari rata-rata latih), dan $\text{shift}_{t,i}$ ($z$-score mean/std jendela input
terhadap statistik latih). Bias awal gate diinisialisasi konservatif ($b_g < 0$) sesuai temuan
Gate 1.

\subsection{Modul 4 (stretch)}
Modul 4 (opsional \textit{stretch}, bukan kontribusi utama; dipindahkan ke Sprint 2 bila tidak
terkejar sebelum Gate 2) menambahkan tiga suku regularisasi: (i) penalti sensitivitas perturbasi
graf $\|f(X,A)-f(X,\tilde{A})\|_1$ dengan $\tilde{A}$ dari perturbasi acak \texttt{candidate\_mask},
(ii) varians loss lintas kelompok berbasis kuantil fitur iklim asal yang tersedia pada $t$ (tanpa
memakai \fpath{regime_labels.csv}), dan (iii) regularisasi prior lunak
$\lambda_{\text{prior}}\|A \odot (1 - M_{\text{bps}})\|_1$.

\subsection{Baseline dan protokol}
Baseline wajib pada split, 5 seed (\texttt{0..4}), dan backbone GRU identik meliputi:
\texttt{naive\_last}, \texttt{seasonal\_naive}, LightGBM varian \texttt{(a)/(b)/(c)} (19/34/45
fitur), Graph WaveNet, MTGNN, \textit{Latent-Regime GNN} disederhanakan ala RCSGNet ($K=3$
state laten dari histori harga), dan \textit{Lag-Graph Tanpa Iklim} disederhanakan ala MillGNN
(bias $-\lambda\ell$). Protokol terkunci \fpath{docs/eval_protocol_v1.md} (\texttt{protocol-v1})
menetapkan 3 \textit{rolling-origin} (\texttt{O1..O3}) dengan embargo 20 hari kerja, 1 holdout
tersegel (\texttt{2026-06-01..2026-10-31}), metrik MAE/RMSE/sMAPE pada Rp/kg (\texttt{y\_mask=True}),
stratifikasi lima rezim (\texttt{dry} $>$ \texttt{extreme\_wet} $>$ \texttt{el\_nino} $>$
\texttt{la\_nina} $>$ \texttt{normal}), \textit{Worst-Regime MAE}, uji Diebold--Mariano
(HAC lag $h-1$, koreksi HLN, Holm--Bonferroni $\alpha=0{,}05$), uji pemulihan jeda sintetis,
ablasi 5 varian, dan perturbasi prior BPS 10--50\%.

\section{Batas klaim}
\begin{table}[H]
\centering
\small
\begin{tabular}{@{}L{2.4cm}L{11.8cm}@{}}
\toprule
Kategori & Rumusan \\
\midrule
AMAN & ``Sepengetahuan kami'', pengondisian distribusi jeda per edge ($\ell=1,\dots,8$ pentad + slot \texttt{null}) pada iklim fisik teramati di provinsi asal ($c_{j,\le t}$) untuk ST-GNN harga cabai 34 provinsi Indonesia; benchmark \textit{availability-safe} terstratifikasi 5 rezim dengan \textit{Worst-Regime MAE} dan uji perturbasi prior graf 10--50\%; perbandingan \textit{head-to-head} melawan rezim laten berbasis harga (ala RCSGNet). \\
BERISIKO & Efektivitas \textit{reliability gate} (Modul 2--3) dan penalti \textit{robustness} (Modul 4, jika dijalankan) hanya boleh dilaporkan sebagai temuan ablasi empiris terkualifikasi, bukan klaim penemuan metode \textit{gating}/invariance baru. \\
JANGAN & Graph adaptif/gating/invariance sebagai novelty; edge = arus pasokan fisik atau kausalitas struktural; ``pertama memakai iklim pada GNN harga pertanian'' atau ``pertama memodelkan rezim pada GNN harga pangan''. \\
\bottomrule
\end{tabular}
\end{table}

\end{document}
